
\chapter{Fisiologia Cardíaca}
\label{cap:12}

\normalfont\rmfamily\upshape
O coração humano bate cerca de cem mil vezes por dia, cerca de 2,5 bilhões
de vezes em uma vida de setenta anos. Em cada ciclo, ele ejeta aproximadamente
70~mililitros de sangue --- cinco litros por minuto em repouso, vinte e
cinco em exercício máximo --- num trabalho contínuo que exige coordenação
fina entre evento elétrico, mecânico e hemodinâmico. Mais do que uma bomba
muscular, o coração é um sistema fisiológico \emph{auto-organizado}: a
frequência cardíaca emerge das propriedades intrínsecas das células do nó
sinoatrial, a contração ordenada depende da propagação elétrica pelo
sistema de condução, e o débito adapta-se à demanda metabólica por vias
autônomas e humorais. Nenhum desses fenômenos é trivial quando observado
de perto; todos se tornaram quantitativamente acessíveis no momento em que
a eletrofisiologia foi inscrita em equações diferenciais.

Este capítulo atravessa esses temas --- eletrofisiologia do cardiomiócito,
acoplamento excitação-contração, geração de força, propagação e arritmias
--- com a intenção de mostrar como eles se compõem em um modelo integrado
do coração como sistema. Antes de tratar dos modelos matemáticos em si,
convém fixar a anatomia e a fisiologia de base.

\section{Anatomia e Fisiologia Cardíaca}
\label{sec:12-anatomia}

\begin{definicaoenv}[Sistema Cardíaco]
O sistema cardíaco é o sistema fisiológico que integra eletrofisiologia,
mecânica, metabolismo e regulação autonômica para bombear sangue de modo
contínuo e adaptativo. Compreende o coração (câmaras, valvas, sistema de
condução), a vasculatura coronária que o irriga, e os mecanismos neuro-humorais
que ajustam frequência e contratilidade às demandas do organismo.
\end{definicaoenv}

O coração humano é dividido em quatro câmaras. Os \emph{átrios} --- direito
(AD) e esquerdo (AE) --- funcionam como câmaras de recepção: coletam o
sangue que retorna (venoso no AD, arterial no AE, este sim após passagem
pulmonar) e o transferem para os ventrículos durante a diástole. Os
\emph{ventrículos} --- direito (VD) e esquerdo (VE) --- são câmaras de
bombeamento: geram pressões suficientes para, respectivamente, propelir o
sangue pela circulação pulmonar e pela circulação sistêmica. O VD opera
contra pressões próximas a 25 mmHg; o VE, contra 120 mmHg --- diferença
essa que torna a parede ventricular esquerda substancialmente mais espessa.

O sangue passa por quatro valvas. As valvas \emph{atrioventriculares}
(tricúspide entre AD e VD; mitral entre AE e VE) impedem refluxo durante
a sístole ventricular. As valvas \emph{semilunares} (pulmonar na saída do
VD; aórtica na saída do VE) fecham-se durante a diástole, impedindo que
o sangue regurgite dos grandes vasos para os ventrículos. O eixo de
fluxo é portanto unidirecional: AD $\to$ VD $\to$ artéria pulmonar
$\to$ pulmões $\to$ AE $\to$ VE $\to$ aorta $\to$ tecidos $\to$ AD.

\subsection{Ciclo Cardíaco}
\label{sub:image12-ciclo}

O ciclo cardíaco alterna dois grandes estados: a \emph{diástole}, em que os
ventrículos relaxam e se enchem, e a \emph{sístole}, em que se contraem e
ejéta o sangue. A diástole pode ser subdividida em enchimento passivo
rápido (quando a valva atrioventricular recém-aberta permite a entrada
do sangue armazenado nos átrios e veias pulmonares), diástase (enchimento
lento) e sístole atrial (contração dos átrios que impulsiona o sangue
final). A sístole ventricular subdivide-se em contração isovolumétrica
(a pressão ventricular sobe com as valvas ainda fechadas, sem mudança de
volume), ejeção (quando a pressão ultrapassa a das artérias, abrindo as
valvas semilunares) e relaxamento isovolumétrico (as valvas fecham, a
pressão cai com volume constante).

A representação canônica dessa dinâmica é o \emph{loop} pressão-volume,
um diagrama fechado que descreve em duas dimensões a relação instantânea
entre pressão intraventricular e volume de sangue. Em frequências
fisiológicas o ciclo completo dura cerca de 0,8~s --- diástole
aproximadamente 0,5~s, sístole 0,3~s ---, com proporções que se ajustam
automaticamente à frequência cardíaca.

\subsection{Sistema de Condução Elétrica}
\label{sub:image12-conducao}

A capacidade do coração de contrair-se de modo coordenado depende de um
sistema de condução dedicado, formado por células musculares
especializadas com automatismo intrínseco e alta velocidade de propagação.

O \emph{nó sinoatrial} (nó SA), localizado no topo do átrio direito, é o
marca-passo natural do coração: as células que o compõem despolarizam
espontaneamente na frequência de 60--100 batimentos por minuto. O
\emph{nó atrioventricular} (nó AV), posicionado na junção entre átrios e
ventrículos, introduz um atraso de condução da ordem de 0,1~s --- crucial
para permitir o enchimento ventricular completo antes da contração. O
\emph{feixe de His} e seus ramos (esquerdo e direito) conduzem rapidamente
o impulso pelos septos; as \emph{fibras de Purkinje} finalmente distribuem
a ativação pelo endocárdio ventricular, garantindo contração quase
simultânea de toda a massa muscular.

A velocidade de condução varia ao longo desse circuito por um fator
superior a cem: 4~m/s nas fibras de Purkinje, 1--1,5~m/s no feixe de His,
0,05~m/s no nó AV. Essa hierarquia é funcional: o atraso nodal AV é a
chave mecânica do acoplamento temporal entre as câmaras.

\subsection{Eletrocardiograma}
\label{sub:image12-ecg}

\begin{exemploenv}[Ondas do ECG e sua origem]
A onda P registra a despolarização dos átrios; o complexo QRS registra a
despolarização dos ventrículos; a onda T registra a repolarização
ventricular (a repolarização atrial é mascarada pelo QRS e geralmente
não é visível). Os intervalos PR (0,12--0,20~s) e QT (0,35--0,44~s)
oferecem leituras quantitativas da condução: o PR mede o tempo entre o
início da ativação atrial e o início da ativação ventricular --- reflete
essencialmente o atraso nodal AV --- e o QT mede a duração total da
ativação ventricular, incluindo platô e repolarização.
\end{exemploenv}

O ECG é, portanto, uma janela indireta para a eletrofisiologia subjacente.
Anormalidades --- arritmias, distúrbios de condução, isquemia --- deixam
assinaturas específicas: onda Q patológica em infarto antigo, supradesnivelamento
do segmento ST em infarto agudo, prolongamento do intervalo QT em canalopatias
congênitas ou adquiridas, complexos QRS alargados em bloqueios de ramo.

\subsection{Circulação Coronariana}
\label{sub:image12-coronarias}

O coração depende de irrigação própria: as artérias coronárias --- direita
(CD), descendente anterior esquerda (DAE) e circunflexa (Cx) --- distribuem
sangue ao miocárdio a partir da raiz da aorta. O fluxo coronariano ocorre
predominantemente durante a diástole, quando o relaxamento miocárdico
permite a perfusão dos vasos intraparietais. A relação pressão-fluxo é
descrita por
\begin{equation}
Q = \frac{\Delta P}{R_{\mathrm{cor}}}
\end{equation}
onde $Q$ é o fluxo (em torno de 250~mL/min em repouso), $\Delta P$ é o
gradiente de pressão entre a aorta e o seio coronário, e $R_{\mathrm{cor}}$
é a resistência do leito coronariano. O miocárdio extrai cerca de 75\% do
oxigênio do sangue arterial --- uma das taxas mais altas do organismo --- e
disso resulta uma \emph{reserva vasodilatadora limitada}: o aumento do
consumo de O$_2$ depende predominantemente de fluxo, e não de maior
extração.

Isquemia surge sempre que a oferta de O$_2$ cai abaixo da demanda.
Causas incluem estenose aterosclerótica fixa (angina estável), ruptura
de placa com trombose (síndrome coronariana aguda), espasmo coronariano
(angina vasoespástica de Prinzmetal) e aumento primário da demanda
(taquicardia, hipertireoidismo). As consequências elétricas ---
correntes de lesão representadas pelo supradesnivelamento ST --- acompanham
o aparecimento dos sintomas dolorosos e definem, em conjunto com
biomarcadores séricos, o diagnóstico clínico contemporâneo.


\section{Eletrofisiologia do Cardiomiócito}
\label{sec:12-eletrofisiologia}

A unidade fundamental do coração é o cardiomiócito: célula alongada,
nucleada centralmente, com mitocôndrias abundantes e retículo sarcoplasmático
extenso. Sua característica elétrica mais notável é o potencial de ação
ventricular, dramaticamente mais longo que o de uma fibra nervosa ---
200 a 400~ms em vez de 1--2~ms --- e organizado em cinco fases que
refletem a ação sequencial de diferentes correntes iônicas.

\subsection{Potencial de Ação Cardíaco}
\label{sub:image12-pa}

O potencial de ação ventricular segue a sequência designada numericamente:

\textbf{Fase 4} é o repouso. A célula mantém $V_m \approx -85$ mV graças
à alta permeabilidade ao K$^+$ via corrente retificadora de entrada
$I_{\mathrm{K1}}$. O potencial aproxima-se do potencial de equilíbrio de
Nernst para o potássio ($E_K \approx -90$~mV).

\textbf{Fase 0} é a despolarização rápida. O estímulo supralimiar abre
canais de Na$^+$ voltagem-dependentes (Nav1.5), produzindo uma corrente
de entrada massiva $I_{\mathrm{Na}}$ que eleva o $V_m$ para $+30$ mV em
menos de 1~ms --- o upstroke do complexo QRS do ECG.

\textbf{Fase 1} é a repolarização inicial. A rápida inativação dos
canais de Na$^+$ e a abertura transitória de corrente de saída de
K$^+$ ($I_{\mathrm{to}}$) produzem uma breve queda de $V_m$ para
cerca de $+10$ mV.

\textbf{Fase 2} é o platô, a característica distintiva do potencial
cardíaco. A entrada de Ca$^{2+}$ via canais tipo L (Cav1.2) --- $I_{\mathrm{Ca},L}$
--- é equilibrada pelas correntes retificadoras retardadas de K$^+$ --- $I_{\mathrm{Kr}}$
(rápida, hERG) e $I_{\mathrm{Ks}}$ (lenta) --- produzindo um platô de
cerca de 200~ms durante o qual o $V_m$ permanece próximo de 0~mV. Esse
platô é a base elétrica do acoplamento excitação-contração discutido na
seção seguinte.

\textbf{Fase 3} é a repolarização final. À medida que os canais de
Ca$^{2+}$ se inativam e as correntes de K$^+$ persistem, o balanço se
inverte --- o $V_m$ cai rapidamente em direção a $E_K$.

\subsection{Equação de Goldman-Hodgkin-Katz em Cardiomiócitos}
\label{sub:image12-ghk}

O potencial de repouso é descrito pela equação de Goldman-Hodgkin-Katz
(GHK) para múltiplos íons:
\begin{equation}
V_m = \frac{RT}{F}\, \ln \!
\left(\frac{P_K[\mathrm{K}^+]_o + P_{\mathrm{Na}}[\mathrm{Na}^+]_o
                + P_{\mathrm{Cl}}[\mathrm{Cl}^-]_i}
            {P_K[\mathrm{K}^+]_i + P_{\mathrm{Na}}[\mathrm{Na}^+]_i
                + P_{\mathrm{Cl}}[\mathrm{Cl}^-]_o}\right)
\end{equation}
onde $P_X$ são as permeabilidades. Em condições fisiológicas, com
permeabilidades relativas $P_K : P_{\mathrm{Na}} : P_{\mathrm{Cl}}
\approx 1 : 0{,}04 : 0{,}45$, a GHK reduz-se aproximadamente à
equação de Nernst para K$^+$, produzindo $V_m \approx -85$~mV. Alterações
do potássio extracelular --- como a hipercalemia, em que $[\mathrm{K}^+]_o$
chega a 7--8~mM em insuficiência renal --- despolarizam o repouso e
contribuem para arritmias.

\subsection{Formalismo de Hodgkin-Huxley}
\label{sub:image12-hh}

A descrição matemática do potencial de ação cardíaco herda do trabalho de
Hodgkin e Huxley (1952) sobre o axônio gigante de lula, formalismo este
que descreve cada canal iônico como gates probabilísticos. Cada corrente
iônica é escrita como:
\begin{equation}
I_{\mathrm{ion}} = g_{\mathrm{ion}} \cdot (V_m - E_{\mathrm{ion}})
\end{equation}
onde a condutância é modulada por \emph{variáveis de gating}:
\begin{equation}
g_{\mathrm{ion}} = \bar{g}_{\mathrm{ion}} \cdot m^p \cdot h^q
\end{equation}
onde $m$ é a probabilidade de o gate de ativação estar aberto, $h$ a de
inativação estar fechado (isto é, a corrente de porta rápida), e
$p$, $q$ são coeficientes estequiométricos (no canal de Na$^+$,
$p = 3$ e $q = 1$).

As variáveis de gating seguem cinética de primeira ordem com parâmetros
dependentes de voltagem:
\begin{equation}
\frac{dm}{dt} = \alpha_m(V)(1 - m) - \beta_m(V) m
            = \frac{m_\infty(V) - m}{\tau_m(V)}
\end{equation}
onde $\alpha_m$ e $\beta_m$ são as taxas de transição, $m_\infty$ a
probabilidade de equilíbrio e $\tau_m$ a constante de tempo. A corrente
macroscópica em cardiomiócitos ventriculares reúne ~10--15 correntes
diferentes, cada uma com sua própria combinação de gates --- totalizando
mais de 30 variáveis de estado nos modelos modernos.

\subsection{Canais Iônicos Principais}
\label{sub:image12-canais}

A corrente de sódio $I_{\mathrm{Na}}$ é mediada pelo canal Nav1.5
codificado pelo gene \textit{SCN5A}. Sua ativação é rápida ($\tau \sim
0{,}1$ ms) e a inativação é bifásica --- rápida ($\tau \sim 1$--$5$~ms)
mais lenta ($\tau \sim 20$--$200$~ms). Mutações de ganho-de-função em
\textit{SCN5A} produzem síndrome do QT longo tipo 3; mutações de
perda-de-função, síndrome de Brugada.

A corrente de cálcio tipo L $I_{\mathrm{Ca},L}$ atravessa o canal Cav1.2
codificado por \textit{CACNA1C}. Sua ativação é mais lenta (com limiar
em torno de -40~mV) e é dupla: gate de voltagem $d$ e gate de
inativação $f$, este último também inativado pela própria elevação do
Ca$^{2+}$ citosólico ($f_{\mathrm{Ca}}$) --- um mecanismo de
auto-regulação fundamental para o acoplamento excitação-contração.

As correntes de potássio são múltiplas. $I_{\mathrm{Kr}}$ (\emph{rapid
delayed rectifier}, hERG/Kv11.1, codificado por \textit{KCNH2}) tem
ativação rápida e inativação importante --- sua retificação interna
significa que ela flui fortemente durante a Fase 2 e desaparece durante
o platô. $I_{\mathrm{Ks}}$ (\emph{slow delayed rectifier}, KCNQ1/KCNE1,
\textit{KCNQ1}/\textit{KCNE1}) ativa muito lentamente. $I_{\mathrm{K1}}$
(retificadora interna, Kir2.1/\textit{KCNJ2}) mantém o repouso.
A corrente $I_{\mathrm{to}}$ (transitória externa, Kv4.3) produz a
Fase 1. Nos cardiomiócitos atriais há ainda $I_{\mathrm{Kur}}$ (ultra-rápida)
e $I_{\mathrm{KACh}}$ (acetilcolina-ativada) --- a primeira é alvo de
drogas contra fibrilação atrial; a segunda responde à estimulação vagal.

\subsection{Período Refratário}
\label{sub:image12-refrat}

A inativação dos canais de Na$^+$ produz o \emph{período refratário
efetivo} (ERP), intervalo durante o qual a célula não pode gerar um novo
potencial de ação, mesmo sob estímulo forte --- é a base do
\emph{período refratário absoluto}, de ~150--200~ms. À medida que os
canais recuperam a disponibilidade, segue-se o \emph{período refratário
relativo}, em que estímulo mais intenso consegue deflagrar um potencial
de ação, mas com amplitude reduzida.

O ERP cardíaco é maior do que o neuronal por mais de cem vezes --- duração
que combina com a duração do ciclo cardíaco. Em frequências mais altas,
o ERP é proporcionalmente mais curto, mas não o suficiente para acomodar
mais ciclos: existe uma frequência máxima intrínseca.

\section{Modelos Computacionais do Cardiomiócito}
\label{sec:12-modelos}

\subsection{Evolução Histórica dos Modelos}
\label{sub:image12-modelos-hist}

O primeiro modelo matemático especificamente cardíaco foi proposto por
Denis Noble em 1962, baseando-se no formalismo de Hodgkin-Huxley para
descrever as fibras de Purkinje (sistema de condução) com quatro correntes
iônicas e seis variáveis de estado. O modelo foi pioneiro em mostrar que o
formalismo HH era geral o suficiente para capturar o automatismo
intrínseco --- ajustando a inclinação da Fase 4, o nó SA produzia
disparos espontâneos.

Beeler e Reuter (1977) introduziram o primeiro modelo ventricular
mamífero, com dinâmica de Ca$^{2+}$ intracelular. Luo e Rudy
elaboraram em duas etapas --- \emph{LR1} em 1991 e \emph{LR2} em 1994 ---
incluindo retículo sarcoplasmático, bombas iônicas e o trocador
Na$^+$/Ca$^{2+}$, conseguindo reproduzir acoplamento
excitação-contração em escala submolecular. Em 2004, ten Tusscher e
Panfilov publicaram um modelo ventricular humano com 12 correntes,
incorporando dados de células epiteliais, M e endocárdicas --- distinção
importante para capturar a gênese da onda T do ECG. Em 2011, O'Hara, Virág,
Varró e Rudy publicaram uma reconstrução baseada em ~150 conjuntos de
dados experimentais humanos, hoje o padrão em simulações clínicas.

\subsection{Modelo Simplificado com Cinco Variáveis}
\label{sub:image12-modelo-simples}

Para finalidades didáticas, modelos com 5--8 variáveis de estado capturam
as propriedades qualitativas do PA ventricular. O sistema de variáveis
$V$, $m$, $h$, $d$, $f$, $x$ (potencial, gate rápido de Na$^+$ em
ativação, gate rápido de Na$^+$ em inativação, ativação e inativação de
Ca$^{2+}$, ativação de K$^+$) integra as EDOs:
\begin{align}
C_m \frac{dV}{dt} &= -(I_{\mathrm{Na}} + I_{\mathrm{Ca},L} + I_K)/C_m \\
\frac{dm}{dt} &= (m_\infty(V) - m)/\tau_m \\
\frac{dh}{dt} &= (h_\infty(V) - h)/\tau_h \\
\frac{dd}{dt} &= (d_\infty(V) - d)/\tau_d \\
\frac{df}{dt} &= (f_\infty(V) - f)/\tau_f \\
\frac{dx}{dt} &= (x_\infty(V) - x)/\tau_x
\end{align}
onde as correntes seguem a forma $I_{\mathrm{ion}} = \bar{g}_{\mathrm{ion}}\, m^p h^q (V - E_{\mathrm{ion}})$.
Estimações das constantes de tempo $\tau_m \approx 0{,}1$ ms, $\tau_h \approx 2$ ms,
$\tau_d \approx 3$ ms, $\tau_f \approx 20$ ms, $\tau_x \approx 100$ ms
capturam as ordens de grandeza das transições iônicas.

A implementação reproduz a forma característica do PA ventricular em
presença de um pulso de corrente supralimiar. O trecho de código a
seguir é representativo:

\begin{lstlisting}[language=Python, caption=PA ventricular simplificado (5 variaveis)]
import numpy as np
from scipy.integrate import odeint

def cardiac_ap(y, t, I_stim):
    V, m, h, d, f, x = y

    # Condutancias maximas e potenciais reversos
    gNa, gCa, gK = 23.0, 0.09, 0.282
    ENa, ECa, EK = 50.0, 60.0, -85.0

    # Estados estacionarios e time constants
    m_inf = 1.0 / (1 + np.exp(-(V + 40) / 5))
    h_inf = 1.0 / (1 + np.exp( (V + 65) / 7))
    d_inf = 1.0 / (1 + np.exp(-(V + 10) / 6.24))
    f_inf = 1.0 / (1 + np.exp( (V + 32) / 8))
    x_inf = 1.0 / (1 + np.exp(-(V + 20) / 15))

    tau_m, tau_h = 0.10, 2.0
    tau_d, tau_f, tau_x = 3.0, 20.0, 100.0

    INa = gNa * m**3 * h * (V - ENa)
    ICa = gCa * d * f * (V - ECa)
    IK  = gK  * x * (V - EK)
    Cm  = 1.0  # uF/cm^2

    dV   = -(INa + ICa + IK) / Cm + I_stim(t)
    dmdt = (m_inf - m) / tau_m
    dhdt = (h_inf - h) / tau_h
    dddt = (d_inf - d) / tau_d
    dfdt = (f_inf - f) / tau_f
    dxdt = (x_inf - x) / tau_x

    return [dV, dmdt, dhdt, dddt, dfdt, dxdt]

# Protocolo: estimulo em t = 5 ms
I_stim = lambda t: 50.0 if (4 <= t <= 5) else 0.0
t = np.linspace(0, 500, 5001)
y0 = [-85.0, 0.01, 0.99, 0.001, 0.99, 0.01]
sol = odeint(cardiac_ap, y0, t, args=(I_stim,))

# APD_90: tempo para o PA cair 90% entre pico e repouso
V = sol[:, 0]
peak_t = t[np.argmax(V)]
APD_90 = t[np.argmax(V < (V[0] + 0.1 * (V.max() - V[0])))] - peak_t
print(f"APD_90 estimado: {APD_90:.1f} ms")
\end{lstlisting}

Modelos progressivamente mais detalhados (Luo-Rudy I, LR2, ten Tusscher,
O'Hara-Rudy) adicionam compartimentos --- citosol, retículo sarcoplasmático,
sub-membrana --- e mecanismos auxiliares --- bombas SERCA, trocador
NCX, correntes de fundo --- reproduzindo complexidade suficiente para
predizer arritmias sob condições clínicas. A escolha do modelo depende
da pergunta: para uma resposta qualitativa sobre o comportamento do
PA, cinco variáveis bastam; para simular perturbações em canais
cardíacos sob tratamento farmacológico, é necessário um modelo de
escala humana.


\section{Acoplamento Excitação-Contração}
\label{sec:12-ec}

A contração muscular cardíaca é desencadeada pelo potencial de ação por
meio de uma cascata de eventos iniciada pelo Ca$^{2+}$ --- fenômeno cuja
compreensão foi peça fundamental na cardiologia do século XX. O
acoplamento excitação-contração (EC) pode ser descrito em dez etapas
essenciais.

\subsection{A Cascata do Cálcio}
\label{sub:image12-cascata}

Quando o potencial de ação despolariza a membrana, canais de Ca$^{2+}$
tipo L (DHPR --- dihydropyridine receptors) abrem, gerando uma pequena
corrente $I_{\mathrm{Ca},L}$ de entrada. Esse \emph{trigger calcium}
--- embora insuficiente para desencadear a contração por si só ---
ativa sensores locais nos receptores de rianodina (RyR2) do retículo
sarcoplasmático (SR), fazendo-os liberar uma quantidade maciça de
Ca$^{2+}$ armazenado. Esse processo, chamado \emph{calcium-induced
calcium release} (CICR), realiza uma amplificação da ordem de dez vezes
entre o sinal de gatilho (DHPR) e o sinal total (RyR).

A concentração citosólica de Ca$^{2+}$ transita assim de cerca de
$0{,}1\,\mu$M em diástole para $1$--$10\,\mu$M em sístole. O Ca$^{2+}$
liga-se à troponina C, alterando a conformação do complexo
troponina-tropomiosina e expondo os sítios de ligação da actina à
miosina. A hidrólise de ATP pela cabeça da miosina gera o \emph{power
stroke}, encurtando o sarcômero.

A recaptação do Ca$^{2+}$ ao final da sístole cabe à bomba SERCA
(sarco/endoplasmic reticulum Ca-ATPase), que retorna o íon ao SR contra
gradiente, e ao trocador NCX (Na-Ca exchanger) da membrana plasmática,
que remove Ca$^{2+}$ do citosol em troca de Na$^+$. O relaxamento
subsequente é passivo em termos de energia química, mas o gradiente de
Ca$^{2+}$ é restaurado a custo de ATP pela SERCA --- e é essa bomba, e não
a miosina, que é o maior consumidor de ATP do ciclo cardíaco.

\subsection{Modelo de Dois Compartimentos para Ca$^{2+}$}
\label{sub:image12-ca-modelo}

O sistema mínimo que captura a regulação do Ca$^{2+}$ intracelular
envolve dois compartimentos --- citosol e SR --- cujas concentrações
evoluem segundo:
\begin{align}
\frac{d[\mathrm{Ca}^{2+}]_i}{dt} &= J_{\mathrm{rel}} - J_{\mathrm{up}} + J_{\mathrm{leak}} + \frac{2I_{\mathrm{Ca},L} - I_{\mathrm{NCX}}}{F V_{\mathrm{cell}}} \\
\frac{d[\mathrm{Ca}^{2+}]_{\mathrm{SR}}}{dt} &= J_{\mathrm{up}} - J_{\mathrm{rel}} - J_{\mathrm{leak}}
\end{align}

Os fluxos principais seguem cinética de Michaelis-Menten para
$J_{\mathrm{up}}$, dependência de RyR para $J_{\mathrm{rel}}$ e diferença
de concentrações para $J_{\mathrm{leak}}$. Em modelos completos (LR2,
O'Hara-Rudy), buffers como troponina, calmodulina e calsequestrin são
adicionados como variáveis suplementares, modificando a concentração
livre que efetivamente regula a contração.

A falha do CICR é a causa proximal de várias cardiopatias. Na
\emph{insuficiência cardíaca sistólica}, a redução da amplitude do
transitório de Ca$^{2+}$ enfraquece a contração. No \emph{dano por
isquemia-reperfusão}, a sobrecarga de Ca$^{2+}$ no início da reperfusão
gera contrações anômalas e arritmias. Nas canalopatias --- mutações em
RyR2 --- o canal libera Ca$^{2+}$ espontaneamente fora do ciclo do PA,
gerando \emph{afterdepolarizations} e taquicardias.

\subsection{Geração de Força e Lei de Frank-Starling}
\label{sub:image12-frank-starling}

A força gerada pelo cardiomiócito depende cooperativamente do Ca$^{2+}$
citosólico com coeficiente de Hill $n \approx 4$--$6$:
\begin{equation}
F = F_{\max} \cdot \frac{[\mathrm{Ca}^{2+}]_i^n}{K_d^n + [\mathrm{Ca}^{2+}]_i^n}\, g(L)
\end{equation}

A função $g(L)$ captura a relação comprimento-tensão --- a \emph{lei de
Frank-Starling}, formulada por Otto Frank (1895) e Ernest Starling (1918):
\begin{exemploenv}[Lei de Frank-Starling]
A força de contração do ventrículo é proporcional ao comprimento inicial
das fibras musculares cardíacas dentro da faixa fisiológica. Quando o
retorno venoso aumenta e o enchimento diastólico é maior, o débito
sistólico aumenta automaticamente --- sem necessidade de regulação
nervosa ou humoral --- para acomodar o volume extra.
\end{exemploenv}

O mecanismo molecular reside na sensibilidade do complexo actina-miosina
ao Ca$^{2+}$ dependente do comprimento do sarcômero. Sarcômeros
distendidos (até o limite de ~$2{,}2\,\mu$m de sobreposição ótima)
apresentam maior sensibilidade e recrutam mais pontes cruzadas.
Sarcômeros demasiadamente distendidos (>2{,}4\,\mu$m) começam a perder
sobreposição e a força cai.

O modelo de dois fatores produz a curva característica: ramo ascendente
até $L_0 \approx 2{,}0$--$2{,}2\,\mu$m, platô em uma faixa estreita,
ramo descendente para comprimentos superiores a $2{,}4\,\mu$m. Em
condições fisiológicas, o coração opera predominantemente na faixa de
comprimento entre 1,9 e $2{,}2\,\mu$m --- sempre no ramo ascendente, com
margem para ajuste.

\begin{lstlisting}[language=Python, caption=Curva forca-comprimento (Frank-Starling)]
import numpy as np
import matplotlib.pyplot as plt

def contraction_force(Ca_i, SL, F_max=100.0, Ca_50=1.0, n_H=4):
    # Forca em funcao de calcio intracelular e comprimento de sarcomero.
    Ca_act = Ca_i**n_H / (Ca_50**n_H + Ca_i**n_H)
    L0 = 2.0
    if   SL < 1.6: length_dep = 0.0
    elif SL < L0:  length_dep = (SL - 1.6) / (L0 - 1.6)
    elif SL < 2.4: length_dep = 1.0
    else:          length_dep = 1.0 - 0.5 * (SL - 2.4)
    return F_max * Ca_act * length_dep

Ca_range = np.linspace(0.05, 10, 100)
SL_values = [1.8, 2.0, 2.2]

plt.figure(figsize=(8, 5))
for SL in SL_values:
    F = [contraction_force(ca, SL) for ca in Ca_range]
    plt.plot(Ca_range, F, lw=2, label=f'SL = {SL} um')
plt.xlabel('[Ca2+]i (uM)')
plt.ylabel('Forca (kPa)')
plt.title('Curva comprimento-tensao (Frank-Starling)')
plt.legend(); plt.grid(alpha=0.3); plt.tight_layout()
plt.savefig('frank_starling.png', dpi=150)
\end{lstlisting}


\section{Propagação e Arritmias}
\label{sec:12-propagacao}

A propagação do potencial de ação pelo miocárdio segue a teoria do cabo
(cardiac cable theory), generalização cardíaca da teoria do cabo neuronal.
A equação fundamental para um tecido unidimensional é:
\begin{equation}
\frac{\partial V}{\partial t} = D \nabla^2 V - \frac{I_{\mathrm{ion}}}{C_m}
\end{equation}
onde $D$ é um coeficiente de difusão efetivo que depende da resistência
do tecido e do gap junction, $\nabla^2$ o operador laplaciano e
$I_{\mathrm{ion}}$ a corrente iônica local do cardiomiócito. A
velocidade de propagação $\theta$ é dada aproximadamente por
\begin{equation}
\theta = \lambda \sqrt{2/\tau}
\end{equation}
onde $\lambda$ é a \emph{constante de comprimento espacial} (raiz da razão
entre condutividade longitudinal e transversal) e $\tau$ é a constante de
tempo da membrana. Em tecido atrial, $\theta$ alcança 0,5--1~m/s;
em ventricular, 0,3--0,5~m/s; nas fibras de Purkinje, 4~m/s.

\subsection{Mecanismos de Arritmias}
\label{sub:image12-arrit}

Três categorias de mecanismos geram arritmias. \textbf{(i)} Anormalidades
de formação do impulso incluem automaticidade anômala (células
auto-ritmícas fora do nó SA) e \emph{atividade deflagrada} --- pós-despolarizações
precoces (EADs, durante a Fase 2 ou 3) ou tardias (DADs, após a
repolarização), frequentemente associadas a mutações de canais ou a
fármacos. \textbf{(ii)} Anormalidades de condução incluem bloqueio
atrioventricular e bloqueios de ramo, com produção de ritmos de escape ou
dessincronização mecânica. \textbf{(iii)} Reentrada é o mecanismo mais
comum de taquiarritmias sustentadas, caracterizado por um circuito de
condução com bloqueio unidirecional --- um tecido ainda em estado
refratário em um sentido, mas capaz de conduzir no outro. Se a frente
de onda ao redor do circuito encontrar tecido excitável ao completar a
volta, persiste indefinidamente.

O critério para reentrada envolve o \emph{comprimento de onda} $\lambda$,
produto entre velocidade de propagação $\theta$ e período refratário
efetivo (ERP):
\begin{equation}
\lambda = \theta \times \mathrm{ERP}
\end{equation}
A condição necessária é que o perímetro do circuito exceda $\lambda$ ---
caso contrário, a frente encontra tecido em refratura ao retornar e é
extinguida.

\subsection{Fibrilação Atrial}
\label{sub:image12-fa}

A fibrilação atrial (FA) é a arritmia sustentada mais comum --- com
prevalência de 1--2\% na população geral, ultrapassando 10\% em maiores
de 80 anos. É caracterizada por ativação atrial caótica, com frequência
de 400--600 batimentos por minuto e perda da contração atrial coordenada.
Mecanismos invocados incluem triggers de focos ectópicos --- em sua
maioria no interior das veias pulmonares --- e substrato estrutural com
remodelamento atrial por fibrose.

O fenômeno de \emph{AF begets AF} (FA gera mais FA) é uma observação
experimental crucial: episódios prolongados de FA reduzem o APD atrial
(por \emph{remodeling} --- inativação de $I_{\mathrm{Ca},L}$ e aumento
de $I_{\mathrm{K1}}$) e promovem estabilidade da própria arritmia. A
intervenção precoce --- cardioversão ou ablação --- é fundamental para
prevenir a cronificação.

Modelos computacionais de FA, especialmente os que incorporam rotores
e múltiplas ondas, têm ajudado a guiar estratégias de ablação:
identificação de alvos alvo (\emph{rotor ablation}), isolamento elétrico
das veias pulmonares, e personalização de protocolos de acordo com
parâmetros observados por imageamento cardíaco.

\subsection{Fibrilação Ventricular e Morte Súbita}
\label{sub:image12-fv}

A fibrilação ventricular (FV) é a causa mais comum de morte súbita
cardíaca. No ECG aparece como atividade irregular sem complexo QRS
organizado; mecanicamente, nenhuma contração coordenada é produzida, o
débito cardíaco cai a zero, e a lesão cerebral instala-se em segundos
a minutos. A desfibrilação elétrica --- única terapia eficaz ---
aplica uma corrente despolarizante massiva (200--360 J em adultos)
através do miocárdio, esgotando os gradientes elétricos e permitindo a
retomada do ritmo sinusal pelos marca-passos naturais.

Os mecanismos da FV envolvem \emph{wavebreak} de ondas espirais --- rotores
inicialmente estáveis fraturam-se em padrões espirais múltiplos,
irregulares, de alta frequência. Modelos computacionais em malha
tridimensional (Fenton-Karma, Aliev-Panfilov para tecidos; modelos celulares
detalhados para nível fino) reproduzem o espalhamento caótico observado.

\subsection{Síndrome do QT Longo}
\label{sub:image12-lqts}

A síndrome do QT longo (LQTS) descreve um conjunto de canalopatias com
prolongamento do intervalo QT no ECG (acima de 440~ms em homens, 460~ms
em mulheres). Três formas genéticas principais são tradicionalmente
distinguidas:

\begin{itemize}
\item \textbf{LQT1} (mutação em KCNQ1): redução de $I_{\mathrm{Ks}}$, a corrente
  de repolarização lenta. Gatilhos típicos são exercício físico,
  especialmente natação.
\item \textbf{LQT2} (mutação em KCNH2/hERG): redução de $I_{\mathrm{Kr}}$, com
  surtos associados a estresse emocional, estímulos auditivos súbitos,
  pós-parto.
\item \textbf{LQT3} (mutação em SCN5A): ganho-de-função em $I_{\mathrm{Na}}$ ---
  persistência da corrente de sódio mantém o platô. Eventos ocorrem em
  repouso ou sono.
\end{itemize}

A LQTS pode também ser adquirida --- diversos medicamentos bloqueiam $I_{\mathrm{Kr}}$
(em particular, antibióticos macrolídeos, antieméticos como ondansetrona,
antipsicóticos como haloperidol) e prolongam o QT. O risco arrítmico --- na
forma de \emph{torsades de pointes} --- motivou protocolos regulatórios
que exigem teste in silico para qualquer nova droga. Os modelos de
O'Hara-Rudy e ten Tusscher são usados rotineiramente para prever
prolongamento QT a partir de perfis de bloqueio iônico.

\begin{exemploenv}[Modulação da velocidade de propagação por drogas]
Bloqueadores de canais de sódio (classe I da classificação Vaughan-Williams)
reduzem o upstroke do PA ($dV/dt|_{max}$), com consequente redução da
velocidade de propagação. Em casos extremos, a propagação falha em
alcançar regiões distantes do miocárdio antes que elas se tornem
refratárias, favorecendo bloqueios de condução. Bloqueadores de canais
de potássio (classe III, como amiodarona) prolongam o APD e, via
aumento do ERP, são geralmente antiarrítmicos. O balanço terapêutico
entre condutância e refratariedade é a base da classificação
antiarrítmica moderna.
\end{exemploenv}

\section{Exercícios}
\label{sec:12-exercicios}

\begin{exercicio}
\textbf{Análise do potencial de ação.} Um PA ventricular tem as seguintes
características: $V_{\mathrm{rest}} = -85$ mV, $V_{\mathrm{pico}} = +30$ mV,
$\mathrm{APD}_{90} = 300$ ms. (a) Calcule a amplitude do PA. (b) Se a
frequência cardíaca é 75 bpm, qual a porcentagem do ciclo em que a célula
está refratária? (c) Como uma droga que bloqueia $I_{\mathrm{Kr}}$ afetaria
o APD?
\end{exercicio}

\begin{exercicio}
\textbf{Equação de Nernst.} A 37°C, calcule o potencial de equilíbrio para
Ca$^{2+}$ dados $[\mathrm{Ca}^{2+}]_o = 2$ mM e $[\mathrm{Ca}^{2+}]_i = 0{,}1\;\mu$M.
Use a relação $E = (58/z) \log_{10}([X]_o/[X]_i)$ mV. O que isso implica
sobre a magnitude da força eletromotriz disponível para entrada de
Ca$^{2+}$ durante o platô?
\end{exercicio}

\begin{exercicio}
\textbf{Acoplamento excitação-contração.} Descreva em prosa o fluxograma
completo do acoplamento EC, identificando os papéis de DHPR, RyR, SERCA
e NCX. Em seguida, explique como uma redução de 50\% em $I_{\mathrm{Ca},L}$
afetaria a contração e por que bloqueadores de canal de Ca$^{2+}$
são efetivos no tratamento de hipertensão.
\end{exercicio}

\begin{exercicio}
\textbf{Comprimento de onda de reentrada.} Considere uma velocidade de
condução $\theta = 0{,}5$ m/s e um ERP $= 250$ ms em tecido atrial
em fibrilação atrial crônica. (a) Calcule o comprimento de onda
$\lambda$. (b) Determine o perímetro mínimo de circuito para sustentar
reentrada. (c) Como o remodelamento atrial (encurtamento de APD por
redução de $I_{\mathrm{Ca},L}$ e aumento de $I_{\mathrm{K1}}$) afeta
$\lambda$ e o limiar para reentrada?
\end{exercicio}

\begin{exercicio}
\textbf{Simulação do PA ventricular.} Implemente o modelo simplificado de
cinco variáveis descrito no capítulo em Python. (a) Reproduza um PA em
resposta a um pulso supra-limiar. (b) Calcule o APD$_{90}$. (c) Varie
$g_{\mathrm{Kr}}$ em ±50\% e meça o efeito sobre APD$_{90}$: o resultado
é consistente com a síndrome do QT longo?
\end{exercicio}

\begin{exercicio}
\textbf{Potencial de equilíbrio múltiplo.} A 37°C, calcule os potenciais
de equilíbrio para K$^+$, Na$^+$ e Cl$^-$ a partir das concentrações da
tabela do capítulo (use a equação GHK simplificada). Em seguida, explique
por que a célula em repouso encontra-se essencialmente em $E_K$, com
$V_m$ ligeiramente mais despolarizado.
\end{exercicio}

\begin{exercicio}
\textbf{Curva de Frank-Starling.} (a) Implemente em Python a função
\texttt{contraction\_force} dada no texto. (b) Plote a força em função
do Ca$^{2+}$ intracelular para três comprimentos de sarcômero (1,9; 2,0;
2,1\,\mu$m). (c) Discuta como alterações da sensibilidade da miosina ao
Ca$^{2+}$ (modeladas por variações do $K_d$) afetam a resposta.
\end{exercicio}

\begin{exercicio}
\textbf{Análise clínica.} Um paciente apresenta FC $= 180$ bpm, ritmo
irregular, ausência de ondas P e complexo QRS estreito. (a) Qual o
diagnóstico mais provável? (b) Explique o mecanismo eletrofisiológico.
(c) Quais opções terapêuticas você consideraria e por quê?
\end{exercicio}

\begin{exercicio}
\textbf{Projeto integrador.} Escolha uma das arritmias discutidas
(fibrilação atrial, fibrilação ventricular, LQT, síndrome de Brugada) e:
(a) formule um modelo computacional simplificado (pode-se usar o modelo
de FitzHugh-Nagumo ou um PA reduzido); (b) reproduza a arritmia por
estimulação apropriada; (c) investigue o efeito de uma droga
antiarrítmica; (d) proponha um protocolo de estimulação de baixa energia
(defibrilação de baixa voltagem) com base na teoria de controle de
rotores. A solução deve ser apresentada com texto, figuras e código
documentado.
\end{exercicio}

\section{Notas Bibliográficas}
\label{sec:12-bibliografia}

{A eletrofisiologia cardíaca é o sistema fisiológico para o qual
existe a correspondência mais estreita entre teoria matemática e
experimento. \textit{Mathematical Physiology} de James Keener e James
Sneyd (2ª edição, Springer) é a referência formal completa para
eletrofisiologia iônica e cardíaca --- cobre Hodgkin-Huxley, modelos
ventriculares e propagação em tecidos. \textit{Physiology of the Heart}
de Arnold Katz (5ª edição, Lippincott) é o texto clássico para os
aspectos fisiológicos com inclinação clínica.

Os artigos seminais que definiram o campo incluem o de Denis Noble
(1962) sobre o primeiro modelo de célula cardíaca em fibras de Purkinje
(adaptando o modelo HH), o de Luo e Rudy (1991) sobre o LR1
(introduzindo dinâmica de Ca$^{2+}$ em cardiomiócito ventricular), o de
ten Tusscher e Panfilov (2006) sobre o modelo ventricular humano com
heterogeneidade transmural, e o de O'Hara, Virág, Varró e Rudy (2011)
sobre o modelo padrão atual baseado em dados humanos extensos.

Para arritmias e clínica, \textit{Cardiac Electrophysiology: From Cell to
Bedside} de Zipes e Jalife (6ª edição, Elsevier) é a referência
abrangente para ligação com a prática. \textit{Cardiovascular Physiology
Concepts} de Richard Klabunde (2ª edição, Lippincott) introduz
eletrofisiologia com clareza de fisiologista. Para simulação computacional,
\textit{Models of Cardiac Tissue Electrophysiology} de Clayton et al.\
\textit{(Progress in Biophysics and Molecular Biology, 2011)} e
\textit{Whole-Heart Modeling} de Trayanova (2011) são revisões
esclarecedoras.

Os recursos computacionais incluem o repositório CellML
(\url{https://models.cellml.org/electrophysiology}), que distribui versões
codificadas dos modelos Luo-Rudy, ten Tusscher, O'Hara-Rudy em SBML e
CellML. O \emph{openCARP} (\url{https://opencarp.org}) é o software
livre padrão para simulações eletrofisiológicas em tecidos 2D/3D. O
PhysioNet (\url{https://physionet.org}) fornece bases de ECG anotadas
essenciais para validação experimental. O \emph{Living Heart Project} da
Dassault Systèmes (\url{https://www.3ds.com/simulia}) é referência em
simulação multi-escala integrando eletromecânica.

}

\begin{thebibliography}{99}
\bibitem{keener} Keener J., Sneyd J. (2009). \textit{Mathematical
Physiology II: Systems Physiology}, 2ª ed. Springer, Nova York.

\bibitem{katz} Katz A.M. (2010). \textit{Physiology of the Heart}, 5ª
ed. Lippincott Williams \& Wilkins, Filadelfia.

\bibitem{noble} Noble D. (1962). A modification of the Hodgkin--Huxley
equations applicable to Purkinje fibre action and pacemaker potentials.
\textit{Journal of Physiology}, 160:317--352.

\bibitem{beeler} Beeler G.W., Reuter H. (1977). Reconstruction of the
action potential of ventricular myocardial fibres. \textit{Journal of
Physiology}, 268:177--210.

\bibitem{luo1} Luo C.H., Rudy Y. (1991). A model of the ventricular
cardiac action potential: depolarization, repolarization, and their
interaction. \textit{Circulation Research}, 68:1501--1526.

\bibitem{luo2} Luo C.H., Rudy Y. (1994). A dynamic model of the cardiac
ventricular action potential: II. Afterdepolarizations, triggered
activity, and potentiation. \textit{Circulation Research}, 74:1097--1113.

\bibitem{tentusscher} ten Tusscher K.H., Panfilov A.V. (2006).
Alternans and spiral breakup in a human ventricular tissue model.
\textit{American Journal of Physiology Heart and Circulatory
Physiology}, 291:H1088--H1100.

\bibitem{ohara} O'Hara T., Virág L., Varró A., Rudy Y. (2011). Simulation
of the undiseased human cardiac ventricular action potential: model
formulation and experimental validation. \textit{PLoS Computational
Biology}, 7:e1002061.

\bibitem{clayton} Clayton R.H., Bernus O., Cherry E.M., Dierckx H.,
Fenton F.H., Mirabella L., Panfilov A.V., Sachse F.B., Seemann G.,
Skarda E. (2011). Models of cardiac tissue electrophysiology: an
overview. \textit{Progress in Biophysics and Molecular Biology},
104:1--21.

\bibitem{trayanova} Trayanova N.A. (2011). Whole-heart modeling:
applications to cardiac arrhythmias and defibrillation. \textit{Circulation
Research}, 109:478--491.

\bibitem{zipes} Zipes D.P., Jalife J. (2013). \textit{Cardiac
Electrophysiology: From Cell to Bedside}, 6ª ed. Elsevier, Filadelfia.

\bibitem{klabunde} Klabunde R.E. (2011). \textit{Cardiovascular
Physiology Concepts}, 2ª ed. Lippincott Williams \& Wilkins,
Filadelfia.

\bibitem{vaughan} Vaughan Williams E.M. (1984). A classification of
antiarrhythmic actions reassessed after a decade of new drugs.
\textit{Journal of Clinical Pharmacology}, 24:129--147.

\bibitem{frank} Frank O. (1895). Zur Dynamik des Herzmuskels.
\textit{Zeitschrift für Biologie}, 32:370--447.

\bibitem{starling} Starling E.H. (1918). \textit{The Linacre Lecture on
the Law of the Heart}. Longmans, Green and Co., Londres.

\bibitem{cellml} Lloyd C.M., Lawson J.R., Hunter P.J., Nielsen P.F.
(2008). The CellML model repository. \textit{Bioinformatics},
24:2122--2123. Disponível em \url{https://models.cellml.org/electrophysiology}.

\bibitem{opencarp} Plank G., Loewe A., Lilienkamp T., Prassl A.,
Fischer C., Pezzuto N., Ravens U., Seemann G. (2021). The openCARP
simulation environment for cardiac electrophysiology. \textit{Computer
Methods and Programs in Biomedicine}, 208:106223. Disponível em
\url{https://opencarp.org}.

\bibitem{physionet} Goldberger A.L., Amaral L.A.N., Glass L., Hausdorff
J.M., Ivanov P.C., Mark R.G., Mietus J.E., Moody G.B., Peng C.K.,
Stanley H.E. (2000). PhysioBank, PhysioToolkit, and PhysioNet:
components of a new research resource for complex physiologic signals.
\textit{Circulation}, 101:e215--e220. Disponível em \url{https://physionet.org}.
\end{thebibliography}
